% neighbor_diagram_v3.tex -- fig:neighbors of the manuscript (2026-10-23, his choice; it
% replaces fig:omega-weight). run.sh copies the pdf to manuscript/figures/.
% Single column: the bounding box is 246 pt wide (the revtex \linewidth), so
% \includegraphics[width=\linewidth] prints every size 1:1.
%
% (a) the point r and three neighbors: a generic one at r + d, with the angle theta_d of d
%     and the angle theta_nb of its rod, both from the x axis, and the angle
%     theta_nb - theta_d between the rod and the line through r and r + d (dashed beyond the
%     neighbor), as in v1; one whose rod points at r (+S, red) and one whose rod lies across
%     the line to r (-S, blue). Colors as in fig:aimed-rods.
% (b) the weight per unit distance, 2 pi d Omega(d), from omega_curve.dat (omega_data.py,
%     which checks: area 1, median 0.682 um, 93.75% within 2 um, peak at 0.538 um). Shaded:
%     the half of the weight within the median radius; dashed: the median and 2 um. As in
%     fig:omega-weight (omega_weight.wls), whose curve color it keeps.
% Build with ./run.sh.
\documentclass[tikz,border=0pt]{standalone}
% neighbor_preamble.tex -- shared by neighbor_diagram_v1/v2/v3.tex.
% Notation as in manuscript/definitions.tex: \r = \mathbf{r}, \bs = \boldsymbol,
% \elln = \ell_{\rm n}, \micron = \upmu\mathrm{m}. The body font is Computer Modern, as in
% revtex (no font package), and standalone sets 10 pt, so \small = 9 pt and
% \footnotesize = 8 pt print at those sizes when the pdf is included at its natural width.
\usepackage{amsmath,amssymb,upgreek}
\usepackage{xcolor}
\usepackage{tikz}
\usetikzlibrary{arrows.meta,calc}
\newcommand{\bs}{\boldsymbol}
\renewcommand{\r}{\mathbf{r}}
\newcommand{\elln}{\ell_{\rm n}}
\newcommand{\micron}{\upmu\mathrm{m}}
% the two sign colors of fig:aimed-rods (aimed_rods_exact.wls): the ends of Mathematica's
% TemperatureMap, cmap[0.97] (raises the variance) and cmap[0.03] (lowers it)
\definecolor{plusS}{rgb}{0.832,0.248,0.178}
\definecolor{minusS}{rgb}{0.226,0.355,0.936}
% the curve color of fig:omega-weight (omega_weight.wls)
\definecolor{curve}{rgb}{0.20,0.40,0.70}
\tikzset{
  >={Stealth[length=4.5pt,width=3.6pt]},
  rod/.style={line width=2.2pt, line cap=round},
  ref/.style={gray!75, line width=0.4pt},
  ang/.style={line width=0.5pt},
  pt/.style={fill=white, draw=black, line width=0.8pt},
  dvec/.style={->, line width=0.7pt},
  panel/.style={font=\small\bfseries},
}
% the point r: white disk with a black ring, as in fig:aimed-rods
\newcommand{\pointr}[1]{\filldraw[pt] (#1) circle (2.3pt);}
% the center of a neighbor rod
\newcommand{\nbdot}[1]{\fill (#1) circle (1.3pt);}

\usepackage{pgfplots}
\pgfplotsset{compat=1.16}
% generated by omega_data.py
\def\dmed{0.6822}%  median radius, um
\def\radmed{1.1314}%  2 pi d Omega at the median
\def\dpeak{0.5383}%  peak of 2 pi d Omega, um
\def\radtwo{0.0625}%  2 pi d Omega at 2 um

\begin{document}
\begin{tikzpicture}[font=\small]
  \useasboundingbox (0,-0.05) rectangle ++(246pt,4.25);
  % ------------------------------------------------------------- (a) geometry
  \node[panel, anchor=north west, inner sep=0pt] at (0,4.20) {(a)};
  \def\thd{40}\def\thn{80}\def\L{2.15}\def\hr{0.62}
  \begin{scope}[shift={(0.95,1.45)}]
    \coordinate (r) at (0,0);
    \coordinate (n) at (\thd:\L);
    % x axis through r, and its parallel through the neighbor
    \draw[ref, ->] (-0.15,0) -- (2.50,0) node[right=-1pt, text=black] {$x$};
    \draw[ref, dashed] (n) -- ++(0.45,0);
    % the line through r and r + d, continued beyond the neighbor
    \draw[ref, dashed] (n) -- ++(\thd:0.78);
    % theta_d at r, theta_nb at the neighbor, both from the x axis
    \draw[ang] (0.72,0) arc[start angle=0, end angle=\thd, radius=0.72];
    \node at (\thd/2:1.00) {$\theta_{\bs d}$};
    \draw[ang] ($(n)+(0.36,0)$) arc[start angle=0, end angle=\thn, radius=0.36];
    \node[anchor=west, inner sep=1pt] at ($(n)+(0.50,0.09)$) {$\theta_\text{nb}$};
    % theta_nb - theta_d at the neighbor, from the line to r to the rod
    \draw[ang] ($(n)+(\thd:0.58)$) arc[start angle=\thd, end angle=\thn, radius=0.58];
    \node[anchor=south west, inner sep=1pt] at ($(n)+(70:0.70)$) {$\theta_\text{nb}\!-\!\theta_{\bs d}$};
    % the generic neighbor
    \draw[dvec, shorten >=1.4pt] (r) -- (n)
      node[pos=0.5, above left=-1pt, inner sep=2pt] {$\bs d$};
    \draw[rod] ($(n)+(\thn:\hr)$) -- ($(n)+(\thn+180:\hr)$);
    \nbdot{n}
    \node[anchor=south east, inner sep=1pt] at ($(n)+(168:0.22)$) {$\r\!+\!\bs d$};
    % a neighbor whose rod points at r, and one whose rod lies across the line to r
    \def\hs{0.36}
    \coordinate (np) at (125:0.98);
    \coordinate (nm) at (290:1.05);
    \draw[ref] (r) -- (np);
    \draw[ref] (r) -- (nm);
    \draw[rod, plusS] ($(np)+(125:\hs)$) -- ($(np)-(125:\hs)$);
    \draw[rod, minusS] ($(nm)+(20:\hs)$) -- ($(nm)-(20:\hs)$);
    \nbdot{np}\nbdot{nm}
    \node[text=plusS, anchor=south west, inner sep=1pt] at ($(np)+(35:0.14)$) {$+S$};
    \node[text=minusS, anchor=west, inner sep=1pt] at ($(nm)+(20:0.40)$) {$-S$};
    \pointr{r}
    \node[anchor=north east, inner sep=1pt] at (-0.06,-0.06) {$\r$};
  \end{scope}
  % ------------------------------------------------------------- (b) the weight
  \node[panel, anchor=north west, inner sep=0pt] at (3.95,4.20) {(b)};
  \begin{axis}[
      at={(4.98cm,0.85cm)}, anchor=south west,
      scale only axis, width=3.52cm, height=2.95cm,
      xmin=0, xmax=3, ymin=0, ymax=1.5,
      xtick={0,1,2,3}, ytick={0,0.5,1,1.5},
      extra x ticks={0.5}, extra x tick labels={$\elln$},
      extra x tick style={tick label style={yshift=0.5pt}},
      yticklabels={$0$,$0.5$,$1$,$1.5$},
      tick label style={font=\footnotesize},
      tick align=inside, tick style={black, line width=0.4pt},
      major tick length=2.5pt,
      axis line style={line width=0.5pt},
      xlabel={$d~(\micron)$}, ylabel={$2\pi d\,\Omega(d)~(\micron^{-1})$},
      label style={font=\small},
      xlabel style={yshift=3pt}, ylabel style={yshift=-3pt},
      clip=false,
    ]
    \addplot[draw=none, fill=curve, fill opacity=0.22] table[x=d, y=rad] {omega_shade.dat}
      \closedcycle;
    \addplot[curve, line width=1.1pt, line cap=round, line join=round]
      table[x=d, y=rad] {omega_curve.dat};
    \draw[gray!70!black, dashed, line width=0.6pt] (axis cs:\dmed,0) -- (axis cs:\dmed,1.5);
    \draw[gray!70!black, dashed, line width=0.6pt] (axis cs:2,0) -- (axis cs:2,1.5);
    \node[anchor=west, font=\footnotesize, inner sep=1.5pt] at (axis cs:\dmed,1.38) {$50\%$};
    \node[anchor=west, font=\footnotesize, inner sep=1.5pt] at (axis cs:2,1.38) {$94\%$};
  \end{axis}
\end{tikzpicture}
\end{document}
